% !TeX root = ../main.tex

\sysusetup{
  keywords  = {超大规模集成电路，布局布线，全局布线，并行计算，动态拥塞感知},
  keywords* = {VLSI, placement and routing, global routing, parallel computing, dynamic congestion awareness},
  keywords** = {Mémoire de bachelor, L'Abstract, Mots-clés},
}

\begin{abstract}
  随着先进工艺节点的持续演进，超大规模集成电路（VLSI）设计规模与复杂度呈指数级增长，传统电子设计自动化（EDA）工具在计算效率与设计质量方面面临严峻挑战。布局与布线作为后端设计流程中的核心环节，其结果直接影响芯片的性能、功耗与面积（PPA）。全局布线问题具有强组合优化特性与复杂约束条件，属于典型的 NP 难问题；同时，传统以串行处理为主的算法难以充分利用现代并行计算平台的算力优势，在处理千万级以上规模电路时面临速度与质量的双重瓶颈。在布局环节，由于设计空间极大，为了提升速度，解析式布局算法需要对大量约束进行一定程度的简化，由此产生了与真实布局中难以处理的细粒度约束之间错配。

  针对上述问题，本文围绕基于并行计算及模块级布局规划优化的 VLSI 高效布局布线算法的课题，从并行布线加速的初始布线工作流和模块级布局规划优化工作流两个子课题展开研究。首先，在初始布线阶段，提出一套应用于全局布线阶段的动态拥塞感知初始布线工作流。该方法在生成树构建与路径选择过程中引入实时更新的拥塞代价模型，使布线决策能够动态反映当前全局资源分布状态，从而缓解传统串行布线算法中静态代价函数导致的局部最优问题。同时针对超大规模电路中全局布线的计算效率问题，本文使用了一种面向并行计算平台的全局布线加速方法。该方法从数据级并行与线网级并行两个层面出发，对布线流程进行重构：一方面，使用GPU加速的矩阵运算流程实现了对全局布线资源的高效并行计算；另一方面，同时考虑所有线网的综合布线代价，实现了全局视角的综合衡量。此外在布局阶段，为了缓解粗粒度解析布局算法精确度不足、细粒度全局搜索方法解空间极大的问题，本文研究了基于多段模拟退火算法的细粒度布局优化算法，通过合理设置领域选择以及多段优化策略，可提升布局设计空间探索的效率，从而高效实现对模块空间分布的细粒度动态调节。

  本文核心贡献可概括为三点，一是相比DGR等可微布线方法，新增动态树更新与训练集缩减机制，避免单轮优化的局部最优问题；二是相比CUGR2等传统组合式布线，引入GPU数据级与线网级双层并行，在质量提升的同时控制耗时增长；三是相比纯解析布局与纯模拟退火路线，提出折中的细粒度优化方案，兼顾调整精度与计算效率，无需改动既有布局流程即可集成。

  实验基于ISPD18/19、ICCAD19等31个公开布线基准与ITC'99的11个布局用例展开。结果显示：布线侧拥塞指标相比DGR降低24.9\%、相比CUGR2降低44.1\%，线长与通孔指标近乎无退化；布局侧平均线长优化达3.86\%，WNS与TNS均满足时序设计要求，密度分布仅有限波动。上述结果验证了所提方法的有效性。

\end{abstract}

\begin{abstract*}
  With the continuous evolution of advanced process nodes, the scale and complexity of very-large-scale integration (VLSI) designs have grown exponentially, posing severe challenges to traditional electronic design automation (EDA) tools in both computational efficiency and design quality. As core stages in backend design, placement and routing directly affect power, performance, and area (PPA). Global routing is a typical NP-hard problem with strong combinatorial optimization characteristics and complex constraints; conventional serial-dominated algorithms struggle to exploit modern parallel computing platforms, facing dual bottlenecks in runtime and quality for ten-million-instance scale circuits. In the placement stage, coarse-grained analytical placement simplifies numerous constraints to improve speed, leading to mismatches with the fine-grained constraints that are difficult to handle in real placement.

  To address these issues, this thesis conducts research on VLSI efficient placement and routing algorithms based on parallel computing and module-level placement optimization, focusing on two sub-topics: a parallel routing-accelerated initial routing workflow and a module-level placement planning optimization workflow. First, in the initial routing stage, a dynamic congestion-aware initial routing workflow for the global routing stage is proposed. This method introduces a real-time updated congestion cost model in the tree construction and path selection process, enabling routing decisions to dynamically reflect the current global resource distribution state, thereby alleviating the local optimum problem caused by static cost functions in traditional serial routing algorithms. Meanwhile, to address the computational efficiency problem of global routing in very-large-scale circuits, this thesis adopts a global routing acceleration method for parallel computing platforms. This method reconstructs the routing flow from two levels: data-level parallelism and net-level parallelism. On one hand, it uses GPU-accelerated matrix operations to achieve efficient parallel computing of global routing resources; on the other hand, it considers the comprehensive routing cost of all nets simultaneously, achieving a global perspective comprehensive measurement. Furthermore, in the placement stage, to alleviate the problems of insufficient accuracy in coarse-grained analytical placement algorithms and the huge solution space of fine-grained global search methods, this thesis studies a fine-grained layout optimization algorithm based on multi-stage simulated annealing. By reasonably setting neighborhood selection and multi-stage optimization strategies, it can improve the efficiency of layout design space exploration, thereby efficiently achieving fine-grained dynamic adjustment of module spatial distribution.

  The core contributions of this thesis are summarized as three points: First, compared with differentiable routing methods like DGR, the proposed workflow adds dynamic tree update and training set pruning mechanisms to avoid local optima in single-round optimization. Second, compared with traditional combinatorial routing methods like CUGR2, it introduces dual-layer parallelism (data-level and net-level) on GPU to control runtime growth while improving quality. Third, compared with pure analytical placement and pure simulated annealing routes, it proposes a compromise fine-grained optimization scheme that balances adjustment precision and computational efficiency, and can be integrated without modifying existing layout flows.

  Experiments are conducted on 31 public routing benchmarks including ISPD18/19, ICCAD19 and others, as well as 11 ITC'99 placement cases. Results show: for the routing side, congestion is reduced by 24.9\% compared with DGR and 44.1\% compared with CUGR2, with no degradation in wirelength and via metrics; for the placement side, average wirelength optimization reaches 3.86\%, WNS and TNS meet timing design requirements, and density distribution only fluctuates limitedly. The above results verify the effectiveness of the proposed methods.
\end{abstract*}

% 法语（第三种语言）摘要，本科不需要就全部注释掉即可，研究生请直接注释
%\begin{abstract**}
%  La longeur de l'abstract français doit faire référence à l'abstract chinois. Le titre de l'abstract français est «ABSTRACT». La première lettre de chaque mot-clé doit être en gros, et les mots-clés doivent être séparer par des virgules et espaces. Les mots-clés de chaque langue doivent être correspondants les uns aux autres.
%\end{abstract**}
